\documentclass[tikz,border=10pt]{standalone}

\usepackage{fontspec}
\usepackage[russian,english]{babel}

\babelfont{rm}{Times New Roman}
\babelfont{sf}{Arial}

\usetikzlibrary{
	arrows.meta,
	positioning,
	fit,
	decorations.pathmorphing,
	calc
}

\begin{document}
	
	\begin{tikzpicture}[
		font=\small,
		>=Latex,
		line cap=round,
		line join=round
		]
		
		% -------------------------------------------------
		% Colors
		% -------------------------------------------------
		\definecolor{pump}{RGB}{70,120,200}
		\definecolor{fiber}{RGB}{220,140,40}
		\definecolor{head}{RGB}{190,70,70}
		\definecolor{problem}{RGB}{210,80,80}
		
		% -------------------------------------------------
		% Pump laser
		% -------------------------------------------------
		\node[
		draw,
		rounded corners=4pt,
		minimum width=2.8cm,
		minimum height=1.5cm,
		fill=pump!8,
		thick
		] (laser) at (0,0)
		{\begin{tabular}{c}
				\textbf{Лазер}\\[-1mm]
		\end{tabular}};
		
		% Small laser output indicator
		\draw[very thick, pump]
		([xshift=0.15cm]laser.east) -- ++(0.45,0);
		
		% -------------------------------------------------
		% Fiber
		% -------------------------------------------------
		\coordinate (fstart) at (2.1,0);
		\coordinate (fend)   at (7.5,0);
		
		% Fiber outer jacket
		\draw[
		very thick,
		fiber!70!black
		]
		(fstart) -- (fend);
		
		% Fiber core
		\draw[
		thick,
		white
		]
		(fstart) -- (fend);
		
		% Fiber label
		\node[
		fill=none,
		inner sep=2pt
		] at ($(fstart)!0.5!(fend)+(0,0.42)$)
		{\textbf{Оптическое волокно доставки}};
		
		% Problem highlight around fiber
		\draw[
		draw=problem,
		very thick,
		rounded corners=6pt,
		dashed
		]
		($(fstart)+(-0.15,-1.4)$)
		rectangle
		($(fend)+(0.15,1.4)$);
		
		\node[
		text=problem,
		fill=white,
		inner sep=2pt
		] at ($(fstart)!0.5!(fend)+(0,-1.95)$)
		{\footnotesize нелинейные эффекты и эффекты распространения};
		
		% -------------------------------------------------
		% Arrow into optical head
		% -------------------------------------------------
		\draw[
		->,
		very thick
		]
		(fend) -- ++(0.7,0);
		
		% -------------------------------------------------
		% Optical head
		% -------------------------------------------------
		\node[
		draw,
		rounded corners=6pt,
		minimum width=3.2cm,
		minimum height=2.0cm,
		fill=head!8,
		thick
		] (head) at (9.6,0)
		{\begin{tabular}{c}
				\textbf{Оптическая голова}\\[-1mm]
				\footnotesize нелинейное преобразование
		\end{tabular}};
		
		% Internal nonlinear element
		\node[
		draw,
		fill=head!15,
		minimum width=1.0cm,
		minimum height=0.75cm,
		rounded corners=2pt
		] (nl) at ([xshift=-0.15cm]head.center)
		{НЛЭ};
		
		% Incoming/outgoing beam inside head
		\draw[
		very thick,
		head!80!black,
		-> 
		]
		([xshift=-1.25cm]head.west) -- (nl.west);
		
		\draw[
		very thick,
		head!80!black,
		-> 
		]
		(nl.east) -- ([xshift=1.15cm]head.east);
		
		% Different-frequency output
		\draw[
		very thick,
		head!70!black,
		-> 
		]
		([xshift=1.15cm]head.east) -- ++(0.9,0.25);
		
		% Problem highlight around head
		\draw[
		draw=problem,
		very thick,
		rounded corners=8pt,
		dashed
		]
		([xshift=-0.25cm,yshift=-1.7cm]head.center)
		rectangle
		([xshift=0.25cm,yshift=1.7cm]head.center);
		
		\node[
		text=problem,
		fill=white,
		inner sep=2pt
		] at ($(head.center)+(0,-1.95)$)
		{\footnotesize нелинейное преобразование};
		
	\end{tikzpicture}
	
\end{document}